\section{Introduction}
\label{sec:intro}

A systems language that manages memory must answer one question before any other:
how does it find the objects that are still in use?
Tracing collectors answer with the program roots, which requires a run-time that can scan stacks and registers,
a dedicated thread or at least a stop-the-world rendezvous, and a representation in which the collector can tell a pointer from an integer.
Reference counting answers with the store itself, which requires a barrier on every pointer write and leaves cycles uncollected.
Both answers are structurally expensive: the first builds machinery that suspends or schedules across the whole program,
the second puts the machinery on the hottest path in the program even when almost nothing is shared.

This paper starts from a different observation.
Most functions in a managed program never create, overwrite, or drop a reference that outlives the function.
That is not an approximation to be measured at run time, it is a property to be proved at compile time.
When the proof holds, the compiler can emit code that frees memory at the end of the owning block and nothing else,
and the whole reclamation run-time can be absent from that function.
extC, a small systems language that compiles to C and checks view lifetimes with a depth-indexed escape analysis,
uses this observation as its organizing principle: the compiler derives the root set statically, decides per translation unit
which of three reclamation tiers to emit, and, where a run-time is emitted, gives every unit of its work a budget.

\noindent\emph{From a mechanism to a contract.}
Being able to switch the run-time off is a qualitative result; it is not yet a promise a latency-sensitive user can sign.
What such a user needs is two numbers: an upper bound on the heap and an upper bound on the pause. Our earlier design
reported a pause \emph{distribution}, which is the wrong object to promise, because a distribution has a tail and the tail is
what the user feels. This paper therefore restates the design as a set of contracts and then proves them:
every step of reclamation work has a work bound that does not depend on the size of the heap or of a region;
the heap obeys a trigger law that ties its peak to the live set; the destruction of a whole frame is deferred through an
$O(1)$ queue so that no single step can pay for the frame; and the admission of an allocation reserves room for the largest
request, so that a single large object cannot overrun the wall between one collection and the next.

\noindent\emph{What the compiler computes.}
The compiler assigns each allocation site to one of five storage substrates: a frame arena, a typed reclamation region,
a container pool, a per-task arena, or a plate from \texttt{std.mem}.
Only the reclamation region needs a collector.
Reference objects that can escape to the function level go there, in a non-moving, non-stop-the-world, incremental
mark-and-sweep over typed slabs whose roots are emitted statically; everything else is freed by dropping its scope.
For the reclamation region, a unit-level proof decides the tier: if no function in the unit needs any reclamation work,
nothing is emitted; if the unit is proved acyclic under a single-level-churn condition, counting is emitted and sweeping is
not; otherwise the mark-and-sweep tier is emitted with counting compiled out.

\noindent\emph{Why a roofline.}
Tuning a reclamation policy has four knobs (the trigger ratio, the batch size of a step, the segment length, and the
reclamation credit) and two hard constraints (memory and pause), and practitioners tune it by folklore.
We show that the trade-off has the shape of a roofline model, in the same sense as the GPU model of attainable
performance: the analogue of arithmetic intensity is the reclamation cost per allocated byte, $\lambda$, and the analogue of
attainable throughput is the allocation rate the mutator reaches once reclamation work is charged to it.
The cost of a round is mark plus sweep, the memory is live set plus a fraction of the trigger, and the keep-up condition
is a lower bound on the trigger ratio. These three laws carve an admissible band out of the plane, and within that band the
optimal policy is the one that pushes both knobs to their walls --- more precisely, to the memory wall and the pause wall.
The model is calibrated with six constants measured on the implementation, and it predicts the measured cost of a round to
within a few percent.

\noindent\emph{Results.}
We evaluate on four workload shapes that isolate promotion, size turnover, pointer chasing, and ring turnover, against nine
arms of C++, Rust, Go, Java and JavaScript built from the same specification, in one interleaved frozen run that also checks
that every arm prints the same answer (Section~\ref{sec:eval}).
extC matches the fastest Rust arm on the promotion shape (39.0 against 39.1~ms), is faster than the best C++ arm on size
turnover (16.0 against 17.9~ms) and on the ring (15.0 against 36.1~ms), and is 1.1\% slower than C++ on pointer chasing while
holding 13\% less peak memory.
No shape shows a single pause above one millisecond; the 99.9th percentile is $10$--$58\,_\mu$s depending on the shape, and the
shape that pauses longest is the one with the largest live set, not the one that allocates most.
Turning reclamation off instead of on is not a free lunch: the allocation-heavy shapes become 1.5--3.7 times slower and
17--59 times larger, because reclamation replaces page faults with reuse.
A mechanism ablation that recompiles the same generated program with one budget changed at a time (ten variants) shows which
constant buys what, and where a choice went against us --- a larger scan chunk, a smaller slab, a padded slot --- the measurement
says so and the design keeps the version that survives it.

\noindent\emph{Contributions.}
\begin{itemize}
\item A description of extC memory management as five substrates and a three-tier compile-time switch, with the exact
  static conditions under which each tier is emitted, and a census of every \texttt{.extc} file in the tree showing that
  68.5\% of the programs that compile need no reclamation run-time at all.
\item A soundness theorem for the full current scheme (no dangling reference) that covers, in addition to the four layers of
  the earlier design, the deferred destruction of regions through the closing queue, the reservation made by admission, and
  the switch between tiers.
\item A bounded-pause theorem: one step of marking walks at most a segment or a wall-clock deadline, one step of sweeping
  releases at most a credit, one step of deferred destruction moves at most a fixed number of slabs, and the sum is bounded
  by constants that do not depend on the heap or on the region.
\item A roofline model of reclamation policy: a cost law, a memory law, a pause law and a keep-up law, an admissibility
  theorem that reduces to one inequality, an optimality theorem that puts the optimum on both walls, and a Pareto frontier;
  six calibrated constants and a validation against measurement.
\item A machine-checked core: 13 Lean~4 modules with no \texttt{sorry} covering region safety, the reclaim rule, the idle
  analysis optimality, and the memory-gap law, with a table mapping each paper theorem to its machine-checked counterpart and
  an explicit statement of what the machine check does not cover.
\item An evaluation from one frozen interleaved run: nine arms over four shapes across six languages, a ten-variant mechanism
  ablation, pause distributions with the total work share, a boundedness experiment that scales the round count by $8\times$,
  code-size measurements, and a census.
\end{itemize}


\noindent\emph{Roadmap.}
Section~\ref{sec:design} describes the design operationally, Section~\ref{sec:formal} states the calculus and the theorems,
Section~\ref{sec:roofline} derives the roofline model, Section~\ref{sec:impl} maps the design onto the compiler, and
Section~\ref{sec:eval} measures it. Section~\ref{sec:limitations} states where the design does not apply.
